\title{Parameter-Efficient Orthogonal Finetuning via Factorization}

\begin{abstract}
The prevalence of large foundation models continues to grow, yet training these models from the ground up remains financially and computationally burdensome. Consequently, there is a pressing need for strategies that allow effective adaptation of pre-trained models to downstream tasks without incurring excessive costs. In this work, we focus on a robust adaptation technique—Orthogonal Finetuning (OFT)—designed for this very purpose. While OFT offers notable generalizability, it still requires a considerable count of trainable parameters due to the inherent high dimensionality of orthogonal matrices. To tackle this challenge, we initiate our study with an analysis of OFT through the lens of information transmission, pinpointing several essential properties for improved parameter efficiency. Drawing inspiration from the efficiency gains achieved by the Cooley-Tukey fast Fourier transform algorithm, we introduce an orthogonal parameterization based on butterfly architectures. Incorporating this parameterization within the OFT framework, we develop a new, parameter-efficient adaptation approach, termed Orthogonal Butterfly (BOFT). This method generalizes OFT—encompassing it as a specific scenario—and establishes a broader framework for orthogonal finetuning. Extensive experimental evaluations are performed, encompassing the adaptation of large vision transformers, prominent language models, and state-of-the-art text-to-image diffusion models across a variety of tasks in both vision and language domains.
\end{abstract}

\section{Introduction}

Breakthroughs such as ChatGPT~\cite{brown2020language,chen2021evaluating} and Stable Diffusion~\cite{rombach2022high} showcase the extraordinary capacity of large foundation models to generalize. The performance enhancements witnessed in these models are typically accompanied by a substantial rise in parameter counts

(\emph{e.g.}, GPT-3 comprises roughly 175B parameters). This considerable model size creates significant hurdles for training foundation models ab initio. As such, enabling widespread advancements within the discipline hinges on methods that adapt these foundational models without starting over from scratch. Specifically, the focus shifts to approaches that effectively transfer pre-trained models to new, downstream applications. Three principal strategies exist for efficient model adaptation: \emph{model finetuning}~\cite{hulora2022,zaken2022bitfit,guo2020parameter,raffel2020exploring,qiu2023controlling,zhang2022adaptive,chavan2023one}, which preserves the model structure and refines only a subset of parameters; \emph{adapter tuning}~\cite{rebuffi2017learning,houlsby2019parameter,pfeiffer2020adapterfusion,lin2020exploring,he2021towards}, which supplements the original model with additional learnable parameters; and \emph{prompt tuning}~\cite{li2021prefix,lester2021power}, where new learnable prefix tokens are attached to model inputs. Of these, model finetuning stands out for its conceptual simplicity, strong performance, and—crucially—for not introducing additional inference-time computational delays.

Model finetuning fundamentally relies on maintaining similarity between the pretrained and finetuned models according to specific metrics, thereby retaining the knowledge acquired during pretraining. Traditional finetuning strategies, for example, often employ a reduced learning rate to limit the difference between the pretrained and finetuned models in a rather heuristic way. Given that weight matrices exhibit a structured organization, a more systematic strategy focuses on conserving the relational structure within these matrices, such as the pairwise angles among neurons. This consideration forms the basis of Orthogonal Finetuning (OFT)~\cite{qiu2023controlling}, a recently developed framework where neurons in each layer undergo transformation by a shared orthogonal matrix, ensuring that the pairwise angular relationships between neurons are rigorously maintained during finetuning. While OFT achieves notable improvements in generalization and convergence, particularly for finetuning text-to-image diffusion models~\cite{qiu2023controlling}, it suffers from the high parameter count required by large orthogonal matrices. To alleviate this issue, OFT adopts a block-diagonal form for these matrices, substantially decreasing parameter count. Nevertheless, this parameter-saving approach introduces constraints: the fixed sparsity pattern of block-diagonal matrices reduces the expressiveness of the transformation and necessitates applying transformations independently across blocks. As a result, these block-diagonal patterns not only diminish the ability to represent a wider class of linear transformations but also introduce possible inductive biases. Additionally, determining an effective sparsity configuration remains an unresolved challenge.

Addressing this challenge hinges on constructing a dense orthogonal matrix in a way that conserves parameters. Although it might initially appear impractical—given that a dense orthogonal matrix of dimension $d$ demands $\mathcal{O}(d^2)$ parameters—we introduce a new strategy that builds such matrices through the composition of several sparse, factorized matrices. This method is predicated on the observation that one can decrease the number of parameters by exchanging spatial requirements for computational time. Since we express the orthogonal matrix as the product of sparse matrices, each training step necessitates repeated multiplications. To contextualize this factorization, we conceptualize generating a dense orthogonal matrix as analogous to transmitting information in a network. Within this framework, forming a dense orthogonal matrix by multiplying sparse matrices is akin to information dissemination on a grid-shaped graph. This transmission-oriented perspective leads us to develop various schemes for sparse matrix factorization that keep the parameter count low, yet retain the expressive capacity required to form dense matrices.

For parameter efficiency within this framework, we look to the butterfly networks found in the Cooley-Tukey fast Fourier transform algorithm~\cite{cooley1965algorithm}, where information exchange between nodes occurs efficiently~\cite{klasing1994broadcasting}. In particular, the butterfly topology inherent to Cooley-Tukey suggests a structure for sparse matrix factorization known as butterfly factorization. Employing this technique enables us to generate a dense matrix of size $d$ by multiplying $\mathcal{O}(\log d)$ sparse matrices, each comprising $\mathcal{O}(d)$ non-zero elements. By constraining each constituent sparse matrix to be orthogonal, we obtain a compact orthogonal parameterization requiring only $\mathcal{O}(d\log d)$ parameters, representing a notable reduction compared to the full parameterization. Building on this efficient structure, we introduce Orthogonal Butterfly~(BOFT), a novel method for parameter-efficient finetuning. BOFT extends OFT as a specific case, presenting a comprehensive framework for orthogonal finetuning. Notably, both the block-diagonal and butterfly architectures exploit matrix sparsity to achieve a reduced parameter set. When compared with the low-rank approach exemplified by LoRA~\cite{hulora2022}, our method illustrates how both low-rank and sparse matrices represent key structured matrix classes~\cite{candes2011robust} that underpin parameter-efficient designs.

In contrast to the block-diagonal configuration leveraged by OFT, which balances between expressive capacity and structural regularization, BOFT incorporates the butterfly pattern to yield a more continuous interpolation between matrices in the full orthogonal group (enhancing expressivity) and the identity matrices (enhancing regularity). This architecture allows for the identification of a narrower, more manageable hypothesis class within the broader set of orthogonal transformations for downstream applications. Given the butterfly structure’s pivotal role in enabling efficient computations in widely adopted linear transformations—including the discrete Fourier and cosine transforms—we posit that integrating this structure in parameter-efficient design inherently fosters an inductive bias, improving generalization and mitigating risks of overfitting. Our rationale is rooted in the butterfly configuration’s ability to reproduce various standard linear transforms—a property unattainable for OFT’s block-diagonal setup. Our key contributions are as follows:

\begin{itemize}[leftmargin=*,nosep]
    \item We examine parameter-efficient orthogonal finetuning from the perspective of information transmission, framing the design of dense, parameter-efficient orthogonal matrices as an information flow challenge within grid-based graphs.
    \item Drawing inspiration from the Cooley-Tukey algorithm’s butterfly architecture, we introduce Orthogonal Butterfly, an orthogonal finetuning technique with heightened parameter efficiency, formulated within the information transmission paradigm.
    \item We articulate several theoretical reasons explaining how BOFT can substantially lower the trainable parameter count while maintaining strong expressivity and stable training. Furthermore, the matrix factorization in BOFT enables a natural weight interpolation mechanism (refer to Figure~\ref{fig:boft_interpolation}).
    \item For the first time, we extend orthogonal finetuning beyond its previous focus on controllable text-to-image generation~\cite{qiu2023controlling}, revealing its broad adaptability as a general finetuning strategy. In particular, we deploy BOFT across a variety of adaptation tasks, from computer vision to natural language processing, where it consistently exceeds current state-of-the-art baselines, demonstrating superior parameter efficiency and enhanced generalization.
\end{itemize}

\section{Related Work}

\textbf{Parameter-efficient finetuning (PEFT)}. As foundation models (\emph{e.g.}, \cite{brown2020language,radford2021learning,rombach2022high,kirillov2023segment}) grow in scale and capability, the challenge of adapting these models to specific downstream applications while minimizing the number of trainable parameters has attracted substantial attention~\cite{treviso2023efficient,ding2023parameter,lialin2023scaling}. Among the broad spectrum of PEFT approaches~\cite{houlsby2019parameter,aghajanyan2020intrinsic,hulora2022,edalati2022krona,wang2022adamix,gheini2021cross,zaken2022bitfit,guo2020parameter,sung2021training,ansell2022composable,lester2021power,li2021prefix,vu2022spot,he2021towards,mao2021unipelt,karimi2021compacter,liu2022few,sung2022lst,chen2022parameter,jia2022visual,chen2022adaptformer,zhang2022neural,jie2023fact,lian2022scaling,luo2023towards}, reparameterization techniques~\cite{aghajanyan2020intrinsic,hulora2022,edalati2022krona} are of particular relevance to our work. LoRA~\cite{hulora2022} enhances pretrained weight matrices by incorporating the product of two low-rank matrices, which has yielded favorable outcomes on natural language benchmarks. While LoRA imposes a uniform rank across all inserted matrices, subsequent efforts~\cite{zhang2022adaptive,valipour2022dylora,zhang2023increlora} have explored adaptive rank adjustments across layers to distribute parameter resources more effectively. The popularity of low-rank reparameterization approaches is evident from their adoption in works such as~\cite{dettmers2023qlora,zi2023delta,chavan2023one} due to their conceptual and practical simplicity. Drawing from the theoretical framework of hyperspherical energy and its link to generalization~\cite{liu2018learning,liu2021orthogonal}, \cite{qiu2023controlling} introduce orthogonal finetuning (OFT), which presents an alternative method for tuning text-to-image diffusion models. In OFT, training is restricted to an orthogonal matrix that transforms neuron activations within the same layer, leading to improvements in both model generalization and training stability relative to LoRA. Nevertheless, OFT generally involves a larger number of trainable parameters compared to LoRA. As such, increasing the parameter efficiency of OFT remains an open and important objective. Furthermore, it is not clear whether OFT extends beyond its current use cases—such as the adaptation of text-to-image diffusion models~\cite{qiu2023controlling}—to other adaptation scenarios. BOFT addresses the parameter inefficiency of OFT by leveraging butterfly factorization. This strategy enables us to broaden the applicability of orthogonal finetuning to general adaptation problems.

\textbf{Butterfly structures}. The radix-2 Cooley-Tukey algorithm achieves efficient computation of the $N$-point discrete Fourier transform by recursively breaking it down into two smaller transforms of size $\frac{N}{2}$, which gives rise to the butterfly structure. These structures can be represented as products of several sparse matrices, collectively known as butterfly matrices. Applications of butterfly matrices include parameterizing orthogonal matrices to eliminate the need for pivoting in Gaussian elimination and to enhance computational efficiency~\cite{parker1995random}, stabilizing the optimization of recurrent neural networks~\cite{jing2017tunable}, and performing kernel approximations~\cite{munkhoeva2018quadrature}. Approaches by~\cite{dao2019learning,dao2019kaleidoscope} develop rapid linear transformations through butterfly parameterizations, while \cite{chen2022pixelated,dao2022monarch} employ butterfly matrices for scalable neural network training. Beyond these, butterfly structures have demonstrated usefulness in areas such as accelerating matrix-vector operations~\cite{michielssen1996multilevel,o2010algorithm}, approximating data-sparse matrices~\cite{li2015butterfly}, and facilitating network transmission~\cite{klasing1994broadcasting,li2003linear,rajasingh2016transmission}. Departing from prior studies, our focus is on utilizing butterfly structures to improve the parameter efficiency of OFT.

\section{An Information Transmission View on Orthogonal Finetuning}

We begin by outlining fundamental concepts in OFT. In order to adapt a pretrained weight matrix, OFT reparameterizes the target weight as the product of a learnable orthogonal matrix and the fixed pretrained weight matrix. Unlike LoRA—which introduces updates to the weights via an additive low-rank matrix—OFT achieves updates through multiplication with an orthogonal matrix. LoRA attains parameter efficiency by constraining updates to a low-rank form, whereas the standard OFT model utilizes a block-diagonal orthogonal matrix structure~\cite{qiu2023controlling}; specifically, reducing the diagonal block size enhances the parameter efficiency in OFT. Figure~\ref{fig:comp_lora} visually contrasts these methods. The rationale behind using orthogonal transformations in OFT lies in their ability to maintain neuron pairwise angles~\cite{liu2018learning,liu2021orthogonal,qiu2023controlling}, thereby conserving the semantic information acquired during pretraining. More precisely, OFT introduces an orthogonal matrix $\bm{R}\in\mathbb{R}^{d\times d}$ for a given pretrained linear weight $\bm{W}^0\in\mathbb{R}^{d\times n}$, updating the forward computation from $\bm{z}=(\bm{W}^0)^\top\bm{x}$ to $\bm{z}=(\bm{R}\bm{W}^0)^\top\bm{x}$, with $\bm{x}\in\mathbb{R}^d$ and $\bm{z}\in\mathbb{R}^n$ denoting the input and output vectors, respectively. OFT achieves zero initialization by initializing $\bm{R}$ as the identity matrix. To preserve orthogonality of $\bm{R}$ during optimization, we adopt Cayley parameterization, as described in \cite{qiu2023controlling,liu2021orthogonal}, so that $\bm{R}=(\bm{I}+\bm{Q})(\bm{I}-\bm{Q})^{-1}$ with $\bm{Q}$ being a skew-symmetric matrix ($\bm{Q}=-\bm{Q}^\top$). For the sake of parameter efficiency, the block-diagonal configuration of $\bm{R}$ is achieved by expressing it as $\text{diag}(\bm{R}_1,\bm{R}_2,\cdots,\bm{R}_r)$, where each $\bm{R}_i\in\mathbb{R}^{b\times b}$ is a smaller orthogonal matrix, and $br=d$. This block-diagonal approach, while reducing parameter count, introduces the constraint that the neuron dimensions (columns in $\bm{W}^0$) are partitioned into $r$ groups, each acted on independently by different orthogonal matrices. Despite the demonstrated practical utility of block-diagonal designs, arbitrarily splitting neuron dimensions by index is largely unmotivated, making the use of dense orthogonal matrices more attractive. Thus, this raises a key question: \emph{Is it possible to devise a dense orthogonal matrix formulation while maintaining parameter efficiency?}

To tackle this issue, we suggest expressing a dense orthogonal matrix as a sequence of several sparse orthogonal matrices. To present a cohesive yet clear framework for analyzing the sparsity patterns in orthogonal matrix decompositions, we reinterpret the task of forming a dense orthogonal matrix within OFT as analogous to an information transmission process. More concretely, constructing a dense matrix $\bm{R}\in\mathbb{R}^{d\times d}$ through the multiplication of $m$ square matrices, namely $\thickmuskip=2mu \medmuskip=2mu \bm{R}=\bm{B}_m\bm{B}_{m-1}\cdots\bm{B}_1$, can be seen as routing information through a $d\times (m+1)$ node lattice, as depicted in Figure~\ref{fig:info_trans}. This perspective arises from the realization that a dense square matrix in $d$ dimensions essentially describes fully connected mappings from $d$ sources to $d$ destinations. For matrix $\bm{R}$, a zero entry $R_{ij}$ signals an inability to transfer information from node $j$ to node $i$, while a non-zero value enables the passage of information. Accordingly, the decomposition of the dense matrix $\bm{R}$ into $\bm{B}_m\bm{B}_{m-1}\cdots\bm{B}_1$ can be interpreted as a series of information exchanges, governed by the interaction patterns determined by each $\bm{B}_i$. The order of information flow mirrors the sequence from $\bm{B}_1$ up to $\bm{B}_n$. For illustration, Figure~\ref{fig:info_trans} displays the factorization $\bm{R}=\bm{B}_5\bm{B}_4\bm{B}_3\bm{B}_2\bm{B}_1$, highlighting both their sparsity patterns and the associated graph structure. This figure represents the cumulative connectivity revealed when the matrix multiplications are unfolded. In this graphical representation, each $\bm{B}_i$ specifies how nodes at layer $i$ are linked to those in layer $i+1$. Explicitly, an element $(j_1,j_2)$ in $\bm{B}_i$ indicates the presence of a directed connection from node $j_2$ at layer $i$ to node $j_1$ at layer $i+1$ (with zeros denoting absence). For the product $\bm{B}_5\bm{B}_4\bm{B}_3\bm{B}_2\bm{B}_1$ to yield a dense result, every node from the initial layer must establish communication with all nodes in the sixth layer. Focusing on a shorter product, such as $\bm{R}=\bm{B}_4\bm{B}_3\bm{B}_2\bm{B}_1$ (corresponding to levels one through five), we observe situations where, for instance, information originating at node $1$ fails to reach node $3$, implying that the $(3,1)$ entry in $\bm{B}_4\bm{B}_3\bm{B}_2\bm{B}_1$ is zero. This relationship between the graph structure and the sparsity layout is consistent for subproducts like $\bm{R}=\bm{B}_3\bm{B}_2\bm{B}_1$ and $\bm{R}=\bm{B}_2\bm{B}_1$. Broadly, achieving a dense matrix $\bm{R}\in\mathbb{R}^{d\times d}$ demands that the sequence $\bm{B}_m,\cdots,\bm{B}_1$ prescribes a collection of directed edges on a $d\times (m+1)$ grid, with each edge linking only nodes from adjacent layers (i.e., adjacent columns), such that any node in the starting layer is able to propagate information to any node in the $(m+1)$-st layer.

Interpreting OFT through the lens of information transmission offers insight into the characteristic sparsity in its factorization matrices. As illustrated in Figure~\ref{fig:blk_diag}, the original OFT employs a block-diagonal design for $\bm{R}$, which, while effective in reducing the parameter count, is incapable of yielding a dense form of $\bm{R}$. Our aim is to synthesize a dense orthogonal matrix from $m$ sparse orthogonal factors, striving to minimize the number of active trainable parameters. The information transmission perspective imposes two principal criteria on the construction: \emph{(i)} \textbf{dense connectivity}, requiring that each initial node maintains at least one connection to every terminal node, and \emph{(ii)} \textbf{minimum free edges}, stipulating the edge count remain as low as the orthogonality condition allows. Importantly, orthogonality introduces nuanced restrictions on permissible edges between layers. For example, each factor $\bm{B}_i$ must be full-rank—which is necessary for orthogonality—necessitating $d$ edges that establish a one-to-one mapping between nodes at consecutive levels, thus enforcing a minimum of $d$ inter-level connections (e.g., 4 in the instance shown in Figure~\ref{fig:info_trans}). These edges, which underpin orthogonality, are not learnable and therefore are excluded from the count of trainable connections, as, for a $d\times d$ orthogonal matrix with $d$ non-zeros, these entries are restricted to values $\pm 1$. Orthogonal constraints diminish the requisite parameters relative to unstructured matrices by creating dependencies among the edges: for instance, with a $2\times 2$ orthogonal matrix, if the $(1,1)$ entry is zero, the $(2,2)$ entry must also be zero—i.e., the omission of one edge enforces the absence of another. Consequently, the set of allowable edge patterns is influenced by orthogonality, sometimes mandating additional edges or reducing them. Observing the sparsity pattern (the non-zeros) in the resulting orthogonal matrix is particularly meaningful in the OFT setting, as only the presence and placement of trainable, nonzero parameters in $\bm{R}$ is relevant—while their numerical values are inconsequential.

A straightforward approach to creating a fully-connected (dense) link between two layers requires $\mathcal{O}(d^2)$ connections (that is, the equivalent of a full $d \times d$ orthogonal matrix), amounting to $d^2-d$ connections—for example, with $d=4$, there are $12$ edges. As illustrated in Figure~\ref{fig:info_trans}, a specific example of feasible matrix decomposition utilizes $10$ edges in total, which is notably fewer than what a single dense orthogonal matrix would require. This perspective provides a basis for examining OFT’s parameter-efficiency using graphical analysis, and the framework facilitates the construction of practical decompositions. We draw upon an established network structure from the Cooley-Tukey algorithm, known as butterfly graphs~\cite{cooley1965algorithm}, which are able to efficiently provide dense connections between $d$ input and $d$ output nodes with only $\mathcal{O}(d\log d)$ edges. For instance, although the configuration in Figure~\ref{fig:info_trans} utilizes $10$ edges for full connectivity, a butterfly network achieves this with just $8$ edges. In what follows, we elaborate on the parameter-efficiency gained from employing the butterfly architecture.

\section{Orthogonal Parameterization by Butterfly Factorization}

The butterfly architecture, initially utilized within the Cooley-Tukey algorithm for fast Fourier transforms, offers a means of accomplishing global interactions: a change localized in the frequency domain manifests as widespread effects in the spatial domain—echoing the nature of our information transfer challenge where every node at the initial layer needs connectivity to all nodes at the terminal layer. Additionally, the butterfly configuration is widely adopted in computer networking~\cite{solihin2015fundamentals,li2011linear} for scalable and efficient communication. Let $k\geq 2$ be a power of $2$; we formally define the \emph{butterfly factor} $\bm{B}^F(k)$ as:

\begin{equation}\label{eq:but_factor}
\begin{aligned}
    \bm{B}^F(k)=\begin{bmatrix}
    \text{diag}(\bm{d}_1) & \text{diag}(\bm{d}_2) \\
    \text{diag}(\bm{d}_3) & \text{diag}(\bm{d}_4)
    \end{bmatrix}\in\mathbb{R}^{k\times k},
\end{aligned}
\end{equation}

where each vector $\bm{d}_i\in\mathbb{R}^{\frac{k}{2}},\forall i$. For dimensions $d=2^N$, we then recursively construct the $d$-dimensional \emph{butterfly matrix} $\bm{B}(d)\in\mathbb{R}^{d\times d}$ via

\begin{equation}
\begin{aligned}
    \!\!\bm{B}(d)=\tilde{\bm{B}}(d,d)\!\cdot\!\begin{bmatrix}
    \bm{B}_1(\frac{d}{2})\!\!\!\!\! & \bm{0} \\
    \bm{0}\!\!\!\!\! &\bm{B}_2(\frac{d}{2})
    \end{bmatrix}= \tilde{\bm{B}}(d,d)\tilde{\bm{B}}(d,\frac{d}{2})\cdots\tilde{\bm{B}}(d,2),
\end{aligned}
\end{equation}

where $\bm{B}_1(\frac{d}{2})$ and $\bm{B}_2(\frac{d}{2})$ denote two butterfly matrices, each of dimension $\frac{d}{2}$. We proceed to introduce the concept of a \emph{butterfly component}, formulated as $\thickmuskip=2mu \medmuskip=2mu \tilde{\bm{B}}(d,k)=\text{diag}(\bm{B}^F_1(k),\cdots,\bm{B}^F_{d/k}(k))$. This structure forms a block-diagonal matrix of size $d\times d$, where each block has dimension $k$, and every $\bm{B}^F_i(k)$ is a butterfly factor as specified in Equation~\ref{eq:but_factor}. We can thus leverage the butterfly matrix to construct a parameterization of an orthogonal matrix. This requires that every factor $\tilde{\bm{B}}(d,k)$ within $\bm{B}(d)$ is itself orthogonal. Focusing first on the block-diagonal case $\tilde{\bm{B}}(d,2)$ (block size $2$), orthogonality of this matrix is ensured via the Cayley transform (equivalent to parameterizing a $2$-dimensional rotation), consistent with methods used in \cite{liu2021orthogonal,qiu2023controlling}. The pattern of non-zero entries in any butterfly component can be interpreted as a permutation of that in $\tilde{\bm{B}}(d,2)$, which means that all butterfly components may be straightforwardly parameterized to maintain orthogonality. Consequently, this approach yields an efficient orthogonal parameterization, where the construction is based on numerous $2\times 2$ orthogonal matrices. Extending the construction per \cite{chen2022pixelated}, we define a block butterfly component $\tilde{\bm{B}}^b(d,k)$ by substituting each entry in $\bm{d}_i$ with a $b\times b$ matrix. To preserve orthogonality in $\tilde{\bm{B}}^b(d,2)$, we parameterize each $2b\times 2b$ block matrix as orthogonal. For butterfly components $\tilde{\bm{B}}^b(d,k)$ with $k>2$, the block-wise permutation of the non-zero pattern in $\tilde{\bm{B}}^b(d,2)$ allows for a comparable orthogonal parameterization. Taken together, these constructions enable the forward propagation in BOFT as follows

\begin{equation}
    \begin{aligned}\nonumber
    \bm{z}=\big(\bm{R}(m,b)\cdot\bm{W}^0\big)^\top\bm{x},~~~~\text{s.t.}~\bigg{\{}\bm{R}(m,b)=\prod_{i=1}^m \tilde{\bm{B}}^b_{(d,i)}~~\&~~\underbrace{\big(\tilde{\bm{B}}^b_{(d,j)}\big)^\top\tilde{\bm{B}}^b_{(d,j)}=\tilde{\bm{B}}^b_{(d,j)}\big(\tilde{\bm{B}}^b_{(d,j)}\big)^\top\!=\bm{I}_d}_{\forall j\in [1, m]}\bigg{\}},
    \end{aligned}
\end{equation}

In this formulation, for simplicity, we let $\tilde{\bm{B}}^b(d,2^{m - i+1})$ be represented as $\tilde{\bm{B}}^b_{(d,i)}$, and $\bm{I}_d$ denotes the identity matrix of dimension $d$. The matrix $\bm{R}(m,b)\in\mathbb{R}^{d\times d}$ is an orthogonal transformation formed as the product of several orthogonal butterfly blocks. For notational convenience, we use BOFT($m$,$b$) as shorthand for BOFT with $\bm{R}(m,\frac{b}{2})$, where $b\geq 2$. When $\thickmuskip=2mu \medmuskip=2mu m=1$, the configuration BOFT($1$,$b$) is equivalent to the block-diagonal OFT~\cite{qiu2023controlling} whose block size is $b$. Similarly, BOFT($1$,$d$) recovers the standard OFT~\cite{qiu2023controlling} with a full, unconstrained orthogonal matrix. When $m=\log \frac{2d}{b}$, the block butterfly matrix $\bm{B}^{\frac{b}{2}}(d)$ serves as $\bm{R}$, producing a dense orthogonal matrix. Generally, BOFT($m$,$b$) introduces $\frac{1}{2}(b-1)dm$ effective trainable parameters for fine-tuning a linear transformation of shape $d\times n$. In the special case where a pure butterfly structure is utilized ($m=\log d, b=2$), BOFT only requires $\mathcal{O}(d\log d)$ parameters. In comparison, a standard OFT employing a full dense orthogonal matrix involves $\mathcal{O}(d^2)$ parameters, while a block-diagonal OFT partitioned into $r$ blocks leverages $\mathcal{O}(bd)$. This contrast demonstrates that a dense orthogonal matrix in original OFT necessitates setting $b=d$ as block size, whereas BOFT enables flexible choices for $b$ to still obtain a dense orthogonal transformation.

\textbf{Identity Initialization in BOFT}. It is customary in fine-tuning frameworks to initialize with the exact weights of a pretrained network, aiming to preserve the performance characteristics of the base model. For instance, LoRA initializes its low-rank matrices as zeros. Analogously, BOFT begins by setting all butterfly matrix components to the identity (i.e., their corresponding skew-symmetric matrices are initialized to zero when using the Cayley parameterization).

\textbf{Multiplicative Dropout in BOFT}. LoRA~\cite{hulora2022} integrates Dropout for its low-rank adaptation parameters as a regularization tool. While conventional Dropout~\cite{srivastava2014dropout} is readily compatible with LoRA, the same is not true for BOFT due to its multiplicative update scheme. To resolve this, we introduce a multiplicative Dropout strategy tailored for BOFT. Since $\bm{R}(m,b)$ is built from $m$ butterfly units, each of which can be permuted into block-diagonal $2b\times 2b$ orthogonal matrices, our approach randomly selects a proportion $p_1$ of these butterfly units and a proportion $p_2$ of their diagonal blocks, subsequently substituting the chosen blocks with identity matrices to achieve the Dropout effect.

\section{Intriguing Insights and Discussions}

\textbf{Expressivity of BOFT}. While the butterfly structure, in conjunction with permutations, is capable of exactly representing several classical fast linear transforms~\cite{dao2019learning,dao2019kaleidoscope} (\emph{e.g.}, fast Fourier transform, Hadamard transform), the ability of our orthogonal butterfly matrix to closely approximate an arbitrary orthogonal matrix has not been comprehensively assessed. To investigate this, we performed a numerical experiment aimed at approximating a randomly generated dense orthogonal matrix~\cite{anderson1987generation} of dimensions $512\times 512$. The outcomes presented in Figure~\ref{fig:para_eff} reflect averages over 10 different random seeds. The vertical axis represents the approximation error, while the horizontal axis corresponds to the number of effective trainable parameters. Each curve with a consistent color represents a BOFT model with a particular block size, where the leftmost marker denotes the error for BOFT($1$,$b$) (\emph{i.e.}, the standard block-diagonal OFT with block size $b$). Generally, BOFT achieves superior parameter efficiency when compared to OFT. For instance, BOFT($9$,$2$) attains greater expressiveness than BOFT($1$,$16$) while utilizing considerably fewer parameters. Models with a smaller $b$ and higher $m$ demonstrate a general trend toward enhanced parameter efficiency; for example, BOFT($6$,$4$) can approximate as well as BOFT($2$,$16$) with a reduced parameter count. Broadly speaking, the butterfly matrix encodes a more constrained structure within the orthogonal group, relative to block-diagonal architectures, endowing BOFT with provably greater expressive capacity than OFT when matched in block size.

\begin{theorem}[Expressivity of BOFT]\label{thm:exp_boft}
    The expressive capacity of BOFT exceeds that of OFT given identical block sizes. To approximate the entire space of orthogonal matrices of size $d$, one can employ a product of butterfly matrices: $\bm{B}_{d-1,1}(d)\bm{B}^\top_{d-1,2}(d)\cdots\bm{B}_{1,1}(d)\bm{B}^\top_{1,2}(d)$, where $\bm{B}_{i,j}(d)$ denotes a butterfly matrix for all $i, j$.
\end{theorem}

Theorem~\ref{thm:exp_boft} further enables a direct extension for BOFT, where the resulting orthogonal matrix generalizes to $\thickmuskip=2mu \medmuskip=2mu \bm{R}^G(m_1,b_1,m_2,b_2,l)=\bm{R}_{l,1}(m_1,b_1)\bm{R}_{l,2}^T(m_2,b_2)\cdots\bm{R}_{1,1}(m_1,b_1)\bm{R}_{1,2}^T(m_2,b_2)$. Here, $\bm{R}_{i,j}^T(m,b)$ stands for the orthogonal matrix constructed within BOFT. With the choices $m_1=m_2=\log d$, $b_1=b_2=2$, and $l=d-1$, $\bm{R}^G(m_1,b_1,m_2,b_2,l)$ is sufficiently expressive to capture the entire orthogonal group, which aligns with the structure known as the kaleidoscope hierarchy~\cite{dao2019kaleidoscope}. It should be highlighted, however, that increased expressiveness does not necessarily correspond to enhanced fine-tuning outcomes; universal expressiveness through full fine-tuning often falls short in practice. Thus, navigating the balance between expressivity and regularization is critical for ensuring the generalization of fine-tuned models. BOFT enlarges the range of available fine-tuning parameters via structural inductive biases, supporting an improved trade-off between expressive power and regularization.

\textbf{Spectral properties}. Orthogonal finetuning outperforms LoRA in terms of spectral preservation, as it maintains the spectral norm of the original weight matrix $\bm{W}^0$ without alteration. To illustrate, consider the singular value decomposition: $\bm{W}^0=\bm{U}\bm{\Sigma}\bm{V}^\top$, where $\bm{U}$ and $\bm{V}$ are orthogonal matrices, and $\bm{\Sigma}$ is a diagonal matrix of singular values. In both OFT and BOFT, finetuning proceeds by left-multiplying with an orthogonal matrix $\bm{R}$, yielding updated weights $\bm{R}\bm{U}\bm{\Sigma}\bm{V}^\top$. This operation preserves the largest singular value (i.e., the spectral norm) of $\bm{W}^0$. Such preservation has been empirically linked to improved training robustness and better generalization~\cite{miyato2018spectral,yoshida2017spectral}. Additional theoretical aspects are discussed in Appendix~\ref{app:sn_property}.

\textbf{Orthogonal finetuning as learning bilinear similarity}. The BOFT method can be interpreted as parameterizing bilinear similarity functions of the form $\bm{w}^0_i\bm{R}\bm{x}$, with $\bm{w}^0_i$ denoting the $i$-th neuron (column) from the pretrained weight matrix $\bm{W}^0$. In this framework, BOFT effectively trains a bilinear similarity matrix $\bm{R}$, but with the significant restriction that $\bm{R}$ must remain orthogonal. This constraint directly connects to themes in distance metric learning~\cite{xing2002distance} and theory surrounding bilinear forms~\cite{roman2005advanced}.

\textbf{Inductive bias and generalization in BOFT}. Within BOFT, the parameterization $\bm{R}(m,b)$ generally constitutes a structured subset of the orthogonal group, imposing restrictions on the hypothesis class and thereby introducing inductive bias. We posit that the inductive bias arising from butterfly factorization enhances generalization, as this structure encapsulates patterns present in several classical linear operations~\cite{dao2019kaleidoscope}, such as the discrete Fourier, sine/cosine, and Hadamard transforms. In addition, the inherent sparsity resulting from matrix factorization in BOFT may further contribute to implicit forms of inductive bias~\cite{gunasekar2017implicit,li2018algorithmic,liu2019neural}.

\textbf{Relation to butterfly-based sparse training}. Numerous studies~\cite{dao2019learning,dao2019kaleidoscope,chen2022pixelated,dao2022monarch} have explored sparse training methods utilizing butterfly parameterizations. These approaches generally entail direct reparameterization of the network's weight matrices using the butterfly structure, with models trained from the ground up. In contrast, \cite{dao2022monarch} investigates a finetuning technique wherein pretrained weights are projected onto modified butterfly matrices, followed by optimization of the projected components for specific downstream applications. BOFT, on the other hand, introduces an alternative finetuning methodology, applying transformations via weight matrices that are shared across layers.

\input{rebuttal-results}

\section{Concluding Remarks and Limitations}

We introduce Orthogonal Butterfly, a broadly applicable, parameter-efficient finetuning approach for foundation models that leverages the butterfly matrix architecture. The central concept behind our enhanced parameter efficiency is the representation of a dense orthogonal matrix as a product of several sparse orthogonal matrices. To facilitate the identification of practical matrix factorizations, we develop a framework grounded in graphical information transmission. Utilizing this paradigm, we find the butterfly structure to be especially effective in accomplishing the targeted sparse orthogonal matrix decompositions. Our experimental results underscore the practical benefits of BOFT across various tasks, including finetuning large language models, large-scale vision networks, and text-to-image generation models. Collectively, these findings establish BOFT’s advantage as a general-purpose finetuning method.

While BOFT yields strong empirical performance, it is not devoid of limitations. Specifically, since the ultimate orthogonal matrix is computed as a product of several orthogonal matrices, the method entails somewhat greater training time overhead than OFT. Addressing the efficiency of BOFT during training thus remains an open area for improvement. Notably, after finetuning, the orthogonal matrices acquired via BOFT can be absorbed into the original model, ensuring that inference speed is not adversely affected. Additionally, it remains unclear whether the butterfly network constitutes the optimal topology for information transmission. Our information transmission perspective suggests potential cross-disciplinary synergies, particularly with computer networking research, which extensively investigates topological efficiency for information flow. We anticipate that exploring alternative network architectures may yield even more efficient constructions for dense orthogonal matrices.